\section{Introduction}\label{sec:intro}

The competition scores an agent bundle on hidden GitHub-issue tasks, with Gemma 4 served by vLLM \citep{kwon2023vllm}.
Third-party analyses put the noise of one hosted score near 1.5 of 58 tasks \citep{afshar2026census,mobassir2026toosmall}.
With one submission per day, a hosted score cannot separate two agents that differ by a few tasks. Local evaluation must
therefore carry development; on a Kaggle kernel with four L4 GPUs our stack finished 35.2 agent runs per hour, serially.
A local evaluator has to answer three questions before its numbers mean anything.

\paragraph{Safety.} An agent run executes model-written shell commands and installs model-edited packages; our first
local setup ran them with no sandbox.

\paragraph{Fidelity.} We call a \emph{rendering trap} a configuration detail that silently changes the token sequence
the model receives while runs still finish and score. The \emph{scorer path} is the scorer's configuration reproduced
locally. In our first setup the served model name contained ``gemma4'', so the agent framework (google-adk 1.36.1) sent
tool results with role \texttt{tool\_responses} instead of \texttt{tool} \citep{adk2026toolresponses}. Prompts rebuilt
from its traces matched the server's token counts on 106 of 114 calls only when rendered with \texttt{tool\_responses}
(0 of 114 with \texttt{tool}). Solve counts hid it: that setup solved 3 and 5 of 24 tasks against 4 and 4 on Kaggle
(Fisher p 0.74 and 0.75).

These traps are public: the role rule \citep{adk2026toolresponses}, thought handling in history
\citep{google2026gemma4prompt,dimondev2026hf44,delta90002026vllm38488}, a newline after thinking text
\citep{gerganov2026hf135}, and the scorer's compaction settings \citep{holbrook2026starter}. We discovered none of them;
we contribute checks that detect them (\cref{sec:conform}) and a measurement of three of them plus a template swap
(\cref{sec:case}).

\paragraph{Inference.} Tasks from one repository succeed and fail together, a small model solves about one task in
five, and a development team looks at many variants. A difference between two arms then needs a paired design, an
analysis that respects repositories, and a rule that is fixed before the outcomes are read.

\paragraph{What this version adds.} This is the extended version of a paper of 3,000 words written for the
competition's paper track \citep{markowitz2026papertrack}. \Cref{sec:contain,sec:conform,sec:compare,sec:case-design}
repeat the text of that paper, and \cref{sec:case,sec:programmes} include the text of its supplement. Everything else is
new here. The short paper describes the kit and one experiment. It leaves out most of what was measured while the kit
was built and used: the throughput of the serving stack, the experiments with the 31B model that came before the work
on the small model, the first experiments with E4B, the runs that compared the workstation with the hosted notebook,
and the three candidate changes of the E4B agent that were taken through a screen and a confirmation. This version
reports all of them, with every run as a released row.

\paragraph{Contributions.}
\begin{itemize}
\item \textbf{A kit.} \kit{} has three modules. \emph{Contain} is a container runner for agent jobs together with the
record of the attacks that seven versions of it failed (\cref{sec:contain}). \emph{Conform} is a set of eight checks
that compare the requests of a local stack with the scorer path (\cref{sec:conform}). \emph{Compare} is a set of paired
tests with repositories as clusters and a pass rule whose behaviour was simulated before use (\cref{sec:compare}).
\item \textbf{A measurement of four prompt divergences} on E4B in a pre-registered experiment with five arms, 231 tasks
and 4,620 jobs (\cref{sec:case}).
\item \textbf{Serving measurements} of three Gemma 4 models on the hardware type of the scorer, including the cost of
serving a LoRA adapter (\cref{sec:serving}).
\item \textbf{The record of \RunsTotal{} further agent runs} in \BlocksTotal{} blocks and \ArmsTotal{} arms:
\RunsThirtyOne{} with the 31B model (\cref{sec:big}) and \RunsEFour{} with E4B (\cref{sec:small,sec:programmes}),
together with \ControlRows{} control gradings of the task pool.
\item \textbf{Three complete screen-and-confirm programmes} for a change of the E4B agent, each with its selection
step, its frozen rule and its outcome (\cref{sec:programmes}), and an analysis of what they say about noise and
selection (\cref{sec:analysis}).
\item \textbf{Data and scripts.} Every table and figure of \cref{sec:serving,sec:big,sec:small,sec:programmes,sec:analysis}
is computed by a released script from released files (\cref{sec:avail}).
\end{itemize}

\paragraph{Main findings.}
Most of the results are null results, and we report them as such.
\begin{itemize}
\item No attack on the runner escaped the container's namespaces or privileges, but nine demonstrated high-severity
vectors came through what the host shared with it, in designs that passed their own escape tests.
\item Four prompt divergences from the scorer path changed the E4B solve rate by less than the experiment could
detect; its simultaneous intervals leave room for a drop of up to \SHSimLo{} points and a gain of up to
\STSimHi{} points.
\item Six configurations of a 31B agent, from a public single agent to a wrapped agent with a shell-only coder, solved
\PubMinSeventy{} to \PubMaxSeventy{} of the same 70 tasks; no pairwise difference was detected, and
\PubBetweenSix{} of the 70 tasks were solved by some configurations and not by others.
\item Each of three E4B candidates led its screen. On new tasks the first two had estimated gains below zero. The third
gained \GainThree{} points, with a one-sided $p$ of \GatePThree{} against a required 0.025: not confirmed.
\item When the unchanged baseline ran again on the same 128 tasks, its total moved by \RepeatHDiff{} points.
\HNever{} of these tasks were solved in none of eight runs and \HAlways{} in all eight, so the remaining
\HBetween{} tasks carry all of the repeat-to-repeat noise.
\end{itemize}

\Cref{sec:setting} gives the setting and the terms. \Cref{sec:serving} reports the serving measurements.
\Cref{sec:contain,sec:conform,sec:compare} describe the three modules. \Cref{sec:big,sec:small} report the experiments
with the two models, \cref{sec:case} the case study, and \cref{sec:programmes} the three candidate programmes.
\Cref{sec:analysis} draws the experiments together; related work, limitations and availability follow.
